\section{The Price of Commitment and the Location of the Risk Boundary}
\label{sec:r9_evidence}
The following experiments change surrender eligibility or first-date investment permissions while preserving the original economic experiment as a reported comparison. All operating classes are reoptimized under the relevant changed primitive. The source programs use the hash-pinned transition constructor, complete later menus, and frozen proposal files of the preceding experiment. No network is retrained, no historical scientific input is replaced, and no settlement-grid result is silently reassigned to a different target.

\subsection{Participation Is Quantitatively Material}
\begin{table}[tbp]
\centering\small
\caption{Participation and Service Procurement at the Original Center}
\label{tab:r9_procurement}
\begin{tabular}{llrrr}
\toprule
Contract & Regime & Duration & Grant $q^*$ & $b_{\min}$ \\
\midrule
$F=0.00,\ m=1$ & Both & 0.124688 & 0.110570 & 0.886775 \\
$F=0.70,\ m=1$ & Adjustment & 0.684865 & 0.775123 & 1.131790 \\
$F=0.70,\ m=1$ & No adjustment & 0.436339 & 0.789425 & 1.809203 \\
$F=0.85,\ m=1$ & Adjustment & 0.890233 & 0.781344 & 0.877685 \\
$F=0.85,\ m=1$ & No adjustment & 0.887898 & 0.816592 & 0.919691 \\
$m=8$ & Adjustment & 0.890233 & 0.781344 & 0.877685 \\
$m=8$ & No adjustment & 0.887898 & 0.816592 & 0.919691 \\
\bottomrule
\end{tabular}
\begin{minipage}{0.98\textwidth}\footnotesize\medskip
The agent chooses the better initial class in each regime. Grants and charges use a separate additive utility numeraire, not managed wealth. Here $C(E)=0$; for positive capacity cost add $C(E)$ to the grant and divide the added amount by duration for the change in $b_{\min}$. The principal obtains the reported flow only during operation. These are not empirical monetary compensating variations.
\end{minipage}
\end{table}

At the original center $(\lambda,d)=(0.125,0.425)$, immediate liquidation pays $G(x_0)=0.0223143551$. The noncancellable adjusted optimum is $-0.7590297258$, whereas the no-adjustment optimum is $-0.7942777218$. Thus the minimum signing grants, before any capacity carrying cost, are $0.7813440809$ and $0.8165920769$. Adjustment reduces the compensation required to undertake the mandate by $0.0352479960$. The signing grant is not a wealth injection into the portfolio problem; its separate-account assumption is essential to these numbers.

Table \ref{tab:r9_procurement} also relates the grant to the amount of service actually delivered. With noncancellable operation, the principal's break-even service flow is lower under adjustment. The calculation does not impose that benefit on the principal: below the displayed threshold, the specified service procurement is not mutually acceptable at the minimum grant. Free surrender after the first interval requires less compensation because it delivers much less operation. It is not a cheaper way of obtaining the same service. For a positive carrying cost $C(E)$, add $C(E)$ to the grant and $C(E)/\mathcal A_{p^*}$ to the break-even service flow.

A concrete positive-cost procurement example illustrates the participation margin. Take available capacity at least $0.85$, use $F=E=0.85$, let $C(E)=0.01$ at that capacity, and let the principal's net service flow be $b=0.90$. With adjustment the minimum grant is $0.7913441$ and the principal's surplus is $0.0098659$. Without adjustment the minimum grant is $0.8265921$ and the principal's surplus is $-0.0274839$. Thus the adjusted mandate is mutually acceptable while the same service-price and capacity terms cannot support the unadjusted mandate. This is a constructive theoretical example in the stated separate-account economy, not an empirical calibration or an optimality claim about the mandate menu.

\subsection{A Finite Surrender Charge and a Nonmonotone Relative Option}
\begin{table}[tbp]
\centering\small
\caption{Surrender Changes Both the Level and the Sign of the Relative Option}
\label{tab:r9_surrender}
\begin{tabular}{rrrrr}
\toprule
$m$ & $F$ & $10^4\Delta^{\rm adj}$ & $10^4\Delta^0$ & $10^4\mathcal O$ \\
\midrule
1 & 0.00 & 5.1294 & 5.1294 & 0.0000 \\
1 & 0.40 & 5.1294 & 5.1294 & 0.0000 \\
1 & 0.70 & 43.0565 & 77.1751 & -34.1186 \\
1 & 0.80 & 1.9674 & -1.9186 & 3.8860 \\
1 & 0.85 & 1.9674 & -2.0933 & 4.0607 \\
2 & 0.00 & 5.4927 & 5.4008 & 0.0919 \\
4 & 0.00 & 5.9646 & 5.9646 & 0.0000 \\
6 & 0.00 & 22.4187 & 23.0450 & -0.6263 \\
7 & 0.00 & 25.6088 & 26.0814 & -0.4726 \\
8 & 0.00 & 1.9674 & -2.0933 & 4.0607 \\
\bottomrule
\end{tabular}
\begin{minipage}{0.98\textwidth}\footnotesize\medskip
Original center, full finite operating menu, four reoptimized class values per row. The first interval remains compulsory. $m=8$ is the original mandate and has no eligible surrender date. Absolute values are monotone in $F$ and $m$ in the required directions; the displayed differences need not be. All 19 center specifications remain in the numerical deposit.
\end{minipage}
\end{table}

At zero charge with the first interval compulsory, all four optimized policies surrender at the first eligible date. The positive and nonpositive values are respectively $-0.0882558999182$ and $-0.0887688376826$ in both regimes. The relative adjustment option is zero. The positive initial action is $(0.8,0,0.8)$; the nonpositive action is $(0.8,0,-1/15)$ in the original finite menu. That menu does not contain zero risky share, and the negative position is not described as risk-free. An independent date-zero-surrender calculation instead gives $G(x_0)$, as predicted by the supersolution argument.

This free-surrender result is not inferred from a few optimized rows alone. At all eight dates, the all-action endpoint check at $d=0.45$ has largest interior residual $-0.0214966453$ for the operating backup of $G$. Affinity in $\lambda$ and nonnegative duration extend the check to the original rectangle. The boundary residual is zero. The continuous benchmark's previously established supersolution inequality has the same economic implication for its own free-stopping problem; this is not a convergence argument between the two models.

At the intermediate charge $F=0.70$, the relative option is \emph{negative}, approximately $-0.00341186$. Adjustment still improves each individual class value. What changes is the class in which it is most valuable. Discounted surrender exposures are $0.248753$ and $0.204167$ in the adjusted positive and nonpositive classes, compared with $0.560607$ and $0.533895$ without adjustment. The two regimes therefore value the termination margin differently. These exposures enter \eqref{eq:r9_delta_derivatives}; they are not a causal decomposition obtained by holding the optimized policy fixed while claiming it remains optimal everywhere.

The charge $0.80$ restores opposite rankings at the center, but not the original uniform conclusion. The positive no-adjustment policy still uses the surrender option, with discounted incidence $0.0689951$. A regional certificate at that charge fails to make its upper class difference negative. At $(\lambda,d)=(0,0.4)$, direct reoptimization gives a positive no-adjustment difference $0.00139238$. The deposit retains this adverse corner as well as the failed certificate. Raising the charge to $0.85$ restores the original values and actions at the center. The uniform result below establishes more than that pointwise observation.

Minimum operating terms provide another way of changing the obligation. With free surrender after two, four, six, or seven intervals, both classes continue to prefer the positive position at the center in both regimes. The relative option is not monotone in the minimum term: for example, it is approximately $-0.0000473$ at seven intervals, and becomes $0.0004061$ under the original eight-interval mandate. Absolute class values have the required monotonicity in the fee and the minimum term. Their differences do not inherit it.

\subsection{An Enforcement--Commitment Frontier on the Original Region}
\begin{table}[tbp]
\centering\small
\caption{Sufficient Enforcement Capacity and the Minimum Operating Term}
\label{tab:r9_enforcement}
\begin{tabular}{rrrr}
\toprule
$m$ & Term (years) & Reachable-state bound & All-node bound \\
\midrule
1 & 0.125 & 0.834860 & 2.530979 \\
2 & 0.250 & 0.764293 & 2.231418 \\
3 & 0.375 & 0.681582 & 1.912037 \\
4 & 0.500 & 0.585751 & 1.570892 \\
5 & 0.625 & 0.474158 & 1.206311 \\
6 & 0.750 & 0.342991 & 0.877726 \\
7 & 0.875 & 0.187395 & 0.492947 \\
8 & 1.000 & 0.000000 & 0.000000 \\
\bottomrule
\end{tabular}
\begin{minipage}{0.98\textwidth}\footnotesize\medskip
Uniform over $(\lambda,d)\in[0,0.25]\times[0.4,0.45]$ and both adjustment regimes, on the original finite menu. Bounds are rounded upward to six decimals after the $10^{-7}$ allowance. A capacity $E$ at least the stated bound implements the original mandatory continuation values with charge $F=E$, provided $E\leq\bar E$. These are sufficient statewise bounds, not the smallest initial-state fees or an optimal contract-design solution.
\end{minipage}
\end{table}

The reachable supports contain 30 nodes after one interval and 735 after seven intervals. Applying Proposition \ref{prop:r9_enforcement} with the two feasible no-adjustment lower policies gives the sufficient regional charges in Table \ref{tab:r9_enforcement}. For surrender available after the first interval, the computed upper bound is $0.8348592806$; the displayed sufficient charge is rounded upward. Using every grid node instead would require a sufficient bound above $2.53$. Reachability therefore matters economically: the guarantee need cover the actual mandate's possible continuation states, not unrelated initial conditions.

These are upper bounds on the statewise enforcement requirement, not estimates of the smallest fee preserving just the initial class ranking. They apply to the original finite action target and the whole rectangle $[0,0.25]\times[0.4,0.45]$. They do not automatically cover the enlarged continuous first-date permission menus studied below. Larger minimum terms reduce the sufficient guarantee capacity. Adjustment weakly reduces the exact statewise requirement by \eqref{eq:r9_capacity_substitute}; the supplement also reports the separately solved center thresholds in both regimes.

\begin{table}[tbp]
\centering\small
\caption{The Original Region with a Finite Surrender Charge}
\label{tab:r9_fee_certificate}
\begin{tabular}{lrrc}
\toprule
Contract & $10^4\Delta^{\rm adj}$ bound & $10^4\Delta^0$ bound & Both signs \\
\midrule
$m=8$ & [0.6869, 3.2473] & [-3.6968, -0.5280] & Yes \\
$F=0.80,\ m=1$ & [0.6869, 3.2473] & [-3.6968, 13.9258] & No \\
$F\in[0.85,0.90],\ m=1$ & [0.6869, 3.2473] & [-3.6968, -0.5280] & Yes \\
\bottomrule
\end{tabular}
\begin{minipage}{0.98\textwidth}\footnotesize\medskip
Each row uses the original rectangle $[0,0.25]\times[0.4,0.45]$, both adjustment regimes, and four law subcells. Chord and count are executed separately and give the same displayed bounds. Every unscaled difference includes $2\times10^{-7}$ arithmetic padding relative to the stored arrays. The fee interval uses a joint fee--benefit--law coefficient bound. The failed $F=0.80$ row and its adverse reoptimized corner are retained.
\end{minipage}
\end{table}

Table \ref{tab:r9_fee_certificate} certifies the original opposite-position comparison for every
\begin{equation}
 (\lambda,d,F)\in[0,0.25]\times[0.4,0.45]\times[0.85,0.90],\qquad m=1.
 \label{eq:r9_fee_box}
\end{equation}
Both streamed upper constructions pass with the stated arithmetic padding. Thus the original conclusion is implemented by a finite, explicitly priced enforcement capacity with a surrender right, not only by formally forbidding exit. The capacity remains a substantive economic restriction: participation requires its carrying cost to be funded, and the free and intermediate-charge outcomes remain different. The same original region, rather than a replacement favorable rectangle, appears in all three rows of the table.

\subsection{Both Investment-Permission Boundaries Move}
\begin{table}[tbp]
\centering\small
\caption{The Two Regime-Specific Permission Boundaries}
\label{tab:r9_permissions}
\begin{tabular}{rrrrr}
\toprule
$\lambda$ & $L^*_{\rm adj}$ & $L^*_0$ & $S^*_{\rm adj}$ & $S^*_0$ \\
\midrule
0.000 & 0.792922 & 0.816033 & 0.506281 & 0.486275 \\
0.125 & 0.789009 & 0.810784 & 0.509481 & 0.490704 \\
0.250 & 0.785115 & 0.805683 & 0.512575 & 0.495063 \\
\bottomrule
\end{tabular}
\begin{minipage}{0.98\textwidth}\footnotesize\medskip
Here $d=0.425$, with no voluntary surrender. The $L$ roots hold $S=0.5$; the $S$ roots hold $L=0.8$. Only the first-date risky-share permission changes. Search is exact over the breakpoint reduction conditional on 121 finite consumption--adjustment pairs and feasible original proposals. Reported roots are floating-point computations, not interval enclosures or continuously optimized consumption and preference-adjustment policies.
\end{minipage}
\end{table}

The breakpoint search reproduces the original finite-menu class values at all nine inspected $(\lambda,d)$ contracts to numerical precision. The maximum over continuous first-date risky shares, conditional on each finite consumption--adjustment pair, therefore does not reveal an omitted within-box improvement in these comparisons. The original optimizers are at both position limits, as shown beside the original decision table. The interior preference-hedging formula remains an interior characterization, not a formula for these constrained optimizers.

Table \ref{tab:r9_permissions} places both switching permissions on the same scale. At the center and with $S=0.5$, the adjusted class changes sign at $L\simeq0.7890095$, and the no-adjustment class at $L\simeq0.8107837$. The baseline $L=0.8$ lies between them. Increasing it to $0.82$ crosses both, so both regimes select a positive position. Holding $L=0.8$ instead, the corresponding short-permission switches are $S\simeq0.5094806$ and $S\simeq0.4907037$. Increasing $S$ from $0.5$ to $0.51$ crosses the adjusted boundary, so both regimes select a nonpositive position. The relative adjustment option remains positive in these two particular permission interventions. Its sign alone therefore does not determine whether adjustment changes the selected position class.

At the baseline center, the long-permission shadow values are approximately $0.0179009$ with adjustment and $0.0194114$ without adjustment. The short-permission shadow values are $0.0200460$ and $0.0225172$. These quantities explain the sensitivity in value units per unit of risky-share permission. The exact interpolation-cell derivatives agree with central finite differences to less than $10^{-10}$ in the reported nine-contract audit. A one- or two-percentage-point change in permissions can be large relative to an initial class-value margin near $2\times10^{-4}$.

\begin{table}[tbp]
\centering\small
\caption{Reoptimized Benefit Crossings under Changed Permissions}
\label{tab:r9_benefit_frontiers}
\begin{tabular}{rrrrr}
\toprule
$\lambda$ & $L$ & $S$ & $d^*_{\rm adj}$ & $d^*_0$ \\
\midrule
0.000 & 0.80 & 0.50 & 0.370085 & 0.558915 \\
0.000 & 0.82 & 0.50 & 0.176862 & 0.386619 \\
0.000 & 0.80 & 0.51 & 0.457404 & 0.636790 \\
0.125 & 0.80 & 0.50 & 0.339574 & 0.515873 \\
0.125 & 0.82 & 0.50 & 0.139861 & 0.334568 \\
0.125 & 0.80 & 0.51 & 0.429666 & 0.597689 \\
0.250 & 0.80 & 0.50 & 0.309087 & 0.473014 \\
0.250 & 0.82 & 0.50 & 0.102936 & 0.282746 \\
0.250 & 0.80 & 0.51 & 0.401952 & 0.558753 \\
\bottomrule
\end{tabular}
\begin{minipage}{0.98\textwidth}\footnotesize\medskip
No voluntary surrender. Later-date policies are reoptimized at every trial benefit on the unchanged full menu. All reported local crossings have positive optimizing duration differences. Local numerical brackets and feature values are deposited; they do not establish global uniqueness or a directed-rounding root enclosure. The original permissions and both adverse first-date extensions are shown at every inspected law probability.
\end{minipage}
\end{table}

The operating-benefit frontiers in Table \ref{tab:r9_benefit_frontiers} are recomputed with continuation policies reoptimized as $d$ changes. They are not obtained by extrapolating a fixed policy's duration. At $\lambda=0.125$, the original local crossing pair is approximately $(0.339574,0.515873)$. With the long permission $0.82$, it becomes $(0.139861,0.334568)$; with the short permission $0.51$, it becomes $(0.429666,0.597689)$. Thus the same central benefit $0.425$ can lie between, above, or below the two boundaries. The movement is consistent with the duration and permission-price identities in \eqref{eq:r9_frontier_derivatives}. The numerical root brackets identify local crossings; global uniqueness in $d$ and directed-rounding enclosures of the roots are not inferred.

The economic content of the example is consequently a priced interaction between preference adjustment, operating commitment, and constrained exposure. Adjustment can reduce the compensation and enforcement required to operate while changing which binding permission is preferred. The sign of that change depends on the surrender margin and on the location of the two permission-sensitive value surfaces. The original cardinal-tilt theorem, the settlement-granularity evidence, and the same-target computational comparisons remain necessary complementary statements; none substitutes for these contract margins.
